% Part: model-theory
% Chapter: basics
% Section: overspill

\documentclass[../../../include/open-logic-section]{subfiles}

\begin{document}

\olfileid{mod}{bas}{ove}
\olsection{અતિપ્રવાહ}

\begin{thm}
\ollabel{overspill} જો વાક્યોના ગણ $\Gamma$ને ગમે તેટલા મોટા સાન્ત
નિદર્શો હોય, તો તેને અનંત નિદર્શ હોય છે.
\end{thm}

\begin{proof}
$\Gamma$ની ભાષામાં ગણનીય સંખ્યામાં નવા અચળો $c_0$, $c_1$, \dots
ઉમેરીને તેનો વિસ્તાર કરો અને ગણ $\Gamma \cup \{c_i \neq c_j :
i \neq j\}$ વિચારો. $\Gamma$ને ગમે તેટલા મોટા સાન્ત નિદર્શો છે એમ
કહેવાનો અર્થ એ છે કે દરેક $m >0$ માટે એવું $n\ge m$ મળે છે કે
$\Gamma$ને ગણાંક~$n$નું નિદર્શ હોય. આમાંથી $\Gamma \cup \{c_i \neq
c_j : i \neq j\}$ સાન્ત રીતે સંતોષ્ય છે તેમ ફલિત થાય છે. સઘનતા મુજબ,
$\Gamma \cup \{c_i \neq c_j : i \neq j\}$ને નિદર્શ $\Struct M$ છે;
તે બધી અસમાનતાઓ $c_i \neq c_j$ સંતોષતું હોવાથી તેનું વ્યાપકક્ષેત્ર
અનંત હોવું જ જોઈએ.
\end{proof}

\begin{prop}
\ollabel{inf-not-fo}
કોઈ પણ પ્રથમ-ક્રમ ભાષાનું એવું કોઈ વાક્ય $!A$ નથી જે
!!a{structure}~$\Struct M$માં ત્યારે અને માત્ર ત્યારે જ સત્ય હોય જ્યારે
એ !!{structure}નું વ્યાપકક્ષેત્ર $\Domain{M}$ અનંત હોય.
\end{prop}

\begin{proof}
જો આવું $!A$ હોત, તો તેનો નિષેધ $\lnot !A$ બધી અને માત્ર સાન્ત
!!{structure}sમાં સત્ય હોત. તેથી તેને ગમે તેટલા મોટા સાન્ત નિદર્શો હોત,
પણ કોઈ અનંત નિદર્શ ન હોત; આ \olref{overspill}નો વિરોધાભાસ છે.
\end{proof}

\end{document}
